\documentclass[exercice,a4]{thedocument}%
\usepackage[french]{babel}%
\usepackage{pstricks}%
\usepackage{textcomp}%
\usepackage{amsmath}%
\startdatakeys
    \source{22GENMATAS1}
    \BO{2026}
    \cycle{4}
    \competencesDuSocle{CA,RA}
    \objectifsApprentissage{C4B2eb}
    \motsCles{trigonométrie,angle,Pythagore,triangle}
\enddatakeys
\begin{document}%
\def\degres{\textsuperscript{\textdegree}}%

\medskip

Dans le triangle ADC rectangle en D, l'angle $\widehat{\text{DCA}}$ mesure $30\degres$.

Le point B est le point du segment [AC] tel que les longueurs DB et CB sont égales.

\emph{La figure ci-dessous n'est pas représentée en vraie grandeur}

\begin{center}
\psset{unit=1.15cm,arrowsize=2pt 3}
\begin{pspicture}(6,4.2)
\pspolygon(0.2,3.5)(0.2,0.3)(5.7,0.3)%ADC
\psline(0.2,0.3)(2.95,1.9)%DB
\uput[ul](0.2,3.5){A} \uput[d](2.95,1.9){B} \uput[dr](5.7,0.3){C} \uput[dl](0.2,0.3){D}
\psarc(5.7,0.3){0.5}{150}{180}\rput(5.,0.46){$30\degres$}
\psframe(0.2,0.3)(0.4,0.5)
\psline[linestyle=dashed]{<->}(0.4,3.7)(5.8,0.6)\rput{-30}(3.1,2.55){10~cm}
\rput{50}(1.6,1.1){/}\rput{50}(1.65,1.14){/}\rput{-20}(4.1,1.24){/}\rput{-20}(4.13,1.19){/}
\end{pspicture}
\end{center}

\begin{enumerate}
\item Calculer la mesure de l'angle $\widehat{\text{DBC}}$.
\item Montrer par le calcul que le segment [AD] mesure $5$~cm.
\item Calculer la longueur DC au millimètre près.
\item Déterminer la nature du triangle ABD.
\end{enumerate}

\medskip

\showthesource
\end{document}
